% Glossary terms from ch08: lines of the form \item[Term] Definition.
\item[Prioritized planning] Planning the agents of a MAPF instance one after another in a priority order, each with a single-agent space-time search that treats the paths of the earlier agents as moving obstacles; fast, sound, but incomplete and suboptimal in general.
\item[Priority order] A permutation $\sigma$ of the agents; $\sigma(1)$ is planned first and never yields to anyone, $\sigma(k)$ yields to everyone.
\item[Reservation table] The sets $R_V$ of (vertex, time) entries and $R_E$ of (move, time) entries of the paths planned so far, together with the parked goal cells; a state or move is blocked if it hits an entry.
\item[Parked goal] A goal cell on which an agent has arrived at time $T$ and stays for ever under stay-at-target; reserved for all $t \ge T$.
\item[Cooperative A*] Silver's name for space-time \astar against a reservation table, with the goal-stay test and a horizon; the low-level search of prioritized planning.
\item[Revised rule] The variant of prioritized planning in which each agent also treats the start cells of all lower-priority agents as static obstacles; needed for the completeness guarantee on well-formed instances.
\item[Well-formed instance] An instance in which every agent has a path from its start to its goal that visits no start or goal of another agent; the revised algorithm succeeds on it for every order.
\item[Wait-then-go path] The witness path of the completeness proof: wait at the start until all higher-priority agents are parked, then follow the endpoint-free path.
\item[HCA*] Hierarchical Cooperative \astar: Cooperative \astar with the exact static distance to the goal, computed by a backward search on the map without agents, as heuristic.
\item[WHCA*] Windowed Hierarchical Cooperative \astar: the reservation table is respected only for the next $w$ steps, a prefix of each plan is executed, and all agents replan with rotated priorities; used for lifelong and online planning, without guarantees.
\item[Longest path first] The ordering rule that plans agents in decreasing length of their independent shortest paths; rarely fails but produces expensive plans.
\item[Random restarts] Running prioritized planning with several random orders and keeping the cheapest plan that succeeds.
\item[Priority-based search] A two-level search over partial priority orders (Ma et al.): the high level branches on a conflicting pair of agents by adding $i \prec j$ or $j \prec i$ and searches depth-first, the low level replans the affected agents in topological order against the paths of their higher-priority agents.
